\Article{217}
Für den Fall, dass $bb-ac$ eine positive Nichtquadratzahl ist, haben wir oben gezeigt, dass sich alle Darstellungen der Zahl $M$ durch die Form $app+2bpq+cqq$ (falls es überhaupt welche gibt) durch eine oder mehrere Formeln der folgenden Art angeben lassen:
\[
p=\tfrac{1}{m}(\mathfrak{A}t+\mathfrak{B}u),\quad q=\tfrac{1}{m}(\mathfrak{C}t+\mathfrak{D}u)
\]
Dabei bezeichnen $\mathfrak{A},\mathfrak{B},\mathfrak{C},\mathfrak{D}$ gegebene ganze Zahlen, $m$ den größten gemeinsamen Teiler der Zahlen $a,2b,c$ und schließlich $t,u$ sämtliche ganzen Zahlen, die der Gleichung $tt-(bb-ac)uu=mm$ genügen. Da alle Werte von $t,u$ sowohl positiv als auch negativ genommen werden können, können wir jede jener Formeln durch \textit{vier} andere ersetzen:
\[
\begin{aligned}
p&=\tfrac{1}{m}(\mathfrak{A}t+\mathfrak{B}u), &\quad q&=\tfrac{1}{m}(\mathfrak{C}t+\mathfrak{D}u)\\
p&=\tfrac{1}{m}(\mathfrak{A}t-\mathfrak{B}u), &q&=\tfrac{1}{m}(\mathfrak{C}t-\mathfrak{D}u)
\end{aligned}
\]
\SourceLeaf{222}
\[
\begin{aligned}
p&=\tfrac{1}{m}(-\mathfrak{A}t+\mathfrak{B}u),&
q&=\tfrac{1}{m}(-\mathfrak{C}t+\mathfrak{D}u)\\
p&=-\tfrac{1}{m}(\mathfrak{A}t+\mathfrak{B}u),&
q&=-\tfrac{1}{m}(\mathfrak{C}t+\mathfrak{D}u)
\end{aligned}
\]
Die Gesamtzahl der Formeln ist jetzt viermal so groß wie vorher; $t$ und $u$ bezeichnen aber nicht mehr alle Zahlen, die der Gleichung $tt-(bb-ac)uu=mm$ genügen, sondern nur die positiven. Jede dieser Formen muss daher einzeln betrachtet und darauf untersucht werden, welche Werte von $t$, $u$ ganzzahlige Werte von $x$, $y$ liefern.

Aus der Formel
\[
p=\tfrac{1}{m}(\mathfrak{A}t+\mathfrak{B}u),\quad
q=\tfrac{1}{m}(\mathfrak{C}t+\mathfrak{D}u)\qquad .\ .\ .\ .\ [1]
\]
folgen diese Werte von $x$, $y$:
\[
x=\frac{\mathfrak{A}t+\mathfrak{B}u+mcd-mbe}{m(bb-ac)},\quad
y=\frac{\mathfrak{C}t+\mathfrak{D}u+mae-mbd}{m(bb-ac)}
\]
Oben haben wir aber gezeigt, dass alle (positiven) Werte von $t$ eine rekurrente Folge $t^0,t',t''$ usw. bilden und dass die zugehörigen Werte von $u$ ebenfalls eine rekurrente Folge $u^0,u',u''$ usw. bilden. Außerdem lässt sich eine Zahl $\rho$ so angeben, dass für jeden gegebenen Modul gilt:
\[
t^\rho\equiv t^0,\quad t^{\rho+1}\equiv t',\quad t^{\rho+2}\equiv t''\ \text{usw.},\quad
u^\rho\equiv u^0,\quad u^{\rho+1}\equiv u'\ \text{usw.}
\]
Als diesen Modul nehmen wir die Zahl $m(bb-ac)$. Der Kürze wegen bezeichnen wir die Werte von $x$, $y$, die sich durch Setzen von $t=t^0$, $u=u^0$ ergeben und denen wir den Index $0$ zuordnen, mit $x^0,y^0$; ebenso die durch Setzen von $t=t'$, $u=u'$ entstehenden Werte mit $x',y'$, denen wir den Index $1$ zuordnen, usw. Man sieht dann ohne Mühe: Sind $x^h,y^h$ ganze Zahlen und ist $\rho$ richtig bestimmt, so sind auch $x^{h+\rho},y^{h+\rho}$, ebenso $x^{h+2\rho},y^{h+2\rho}$ und allgemein $x^{h+k\rho},y^{h+k\rho}$ ganze Zahlen. Ist dagegen $x^h$ oder $y^h$ gebrochen, so ist auch $x^{h+k\rho}$ beziehungsweise $y^{h+k\rho}$ gebrochen. Daraus folgt leicht: Berechnet man die Werte von $x$, $y$ mit den Indizes $0,1,2\,.\,.\,.\,.\,\rho-1$, und sind für keinen dieser Indizes \textit{sowohl} $x$ \textit{als auch} $y$ ganzzahlig, so gibt es überhaupt keinen Index, für den sowohl $x$ als auch $y$ ganze Werte annehmen. In diesem Fall lassen sich aus Formel [1] keine ganzzahligen Werte von $x$, $y$ ableiten. Gibt es aber unter jenen Indizes einige, etwa $\mu,\mu',\mu''$ usw., denen ganzzahlige Werte von $x$, $y$ entsprechen, so sind sämtliche ganzzahligen Werte von $x$, $y$, die sich aus Formel [1] gewinnen lassen, diejenigen, deren Indizes in einer der Formeln $\mu+k\rho,\ \mu'+k\rho,$
\SourceLeaf{223}
$\mu''+k\rho$ usw. enthalten sind. Dabei durchläuft $k$ alle positiven ganzen Zahlen, einschließlich der Null.

Die übrigen Formeln, in denen die Werte von $p$, $q$ enthalten sind, müssen genauso behandelt werden. Sollte keine dieser Formeln ganzzahlige Werte von $x$, $y$ liefern, so wäre die vorgelegte Gleichung überhaupt nicht in ganzen Zahlen lösbar. Ist sie aber tatsächlich lösbar, so lassen sich alle ganzzahligen Lösungen nach den oben angegebenen Regeln darstellen.

\Needspace{8\baselineskip}
\Article{218}
Ist $bb-ac$ eine Quadratzahl und $M=0$, so sind alle Werte von $p$, $q$ in zwei Formeln der Art $p=\mathfrak{A}z$, $q=\mathfrak{B}z$; $p=\mathfrak{A}'z$, $q=\mathfrak{B}'z$ enthalten. Dabei durchläuft $z$ alle ganzen Zahlen, während $\mathfrak{A},\mathfrak{B},\mathfrak{A}',\mathfrak{B}'$ gegebene ganze Zahlen sind, von denen die erste und zweite sowie die dritte und vierte jeweils teilerfremd sind (Artikel 212). Alle ganzzahligen Werte von $x$, $y$, die aus der ersten Formel entstehen, sind also in Formel [1] enthalten:
\[
x=\frac{\mathfrak{A}z+cd-be}{bb-ac},\quad
y=\frac{\mathfrak{B}z+ae-bd}{bb-ac}
\]
Alle übrigen, die aus der zweiten Formel entstehen, sind in dieser Formel [2] enthalten:
\[
x=\frac{\mathfrak{A}'z+cd-be}{bb-ac},\quad
y=\frac{\mathfrak{B}'z+ae-bd}{bb-ac}
\]
Da aber beide Formeln auch gebrochene Werte liefern können (außer wenn $bb-ac=1$), müssen wir in jeder Formel diejenigen Werte von $z$, die sowohl $x$ als auch $y$ ganzzahlig machen, von den übrigen trennen. Es genügt jedoch, nur die erste Formel zu betrachten, da auf die andere genau dasselbe Verfahren anzuwenden ist.

Da $\mathfrak{A}$, $\mathfrak{B}$ teilerfremd sind, können wir zwei Zahlen $\mathfrak{a}$, $\mathfrak{b}$ so bestimmen, dass $\mathfrak{a}\mathfrak{A}+\mathfrak{b}\mathfrak{B}=1$ gilt. Dann ist
\[
(\mathfrak{a}x+\mathfrak{b}y)(bb-ac)=z+\mathfrak{a}(cd-be)+\mathfrak{b}(ae-bd)
\]
Hieraus ist sofort klar, dass alle Werte von $z$, die ganzzahlige Werte von $x$, $y$ erzeugen können, notwendig zur Zahl $\mathfrak{a}(be-cd)+\mathfrak{b}(bd-ae)$ modulo $bb-ac$ kongruent sein müssen, also in der Formel $(bb-ac)z'+\mathfrak{a}(be-cd)+\mathfrak{b}(bd-ae)$ enthalten sind, wobei $z'$ sämtliche ganzen Zahlen durchläuft. Damit erhalten wir leicht anstelle der Formel [1] die folgende:
\SourceLeaf{224}
\[
\begin{aligned}
x&=\mathfrak{A}z'+\mathfrak{b}\times\frac{\mathfrak{A}(bd-ae)-\mathfrak{B}(be-cd)}{bb-ac}\\
y&=\mathfrak{B}z'-\mathfrak{a}\times\frac{\mathfrak{A}(bd-ae)-\mathfrak{B}(be-cd)}{bb-ac}
\end{aligned}
\]
Sie liefert offenbar entweder für alle Werte von $z'$ oder für keinen ganzzahlige Werte von $x,y$. Der erste Fall tritt immer dann ein, wenn $\mathfrak{A}(bd-ae)$ und $\mathfrak{B}(be-cd)$ modulo $bb-ac$ kongruent sind, der zweite, wenn sie inkongruent sind. --- Formel [2] ist genauso zu behandeln; ihre ganzzahligen Lösungen (falls sie welche liefern kann) sind von den übrigen zu trennen.

\Needspace{8\baselineskip}
\Article{219}
Ist $bb-ac=0$, so lässt sich die Form $axx+2bxy+cyy$ als $m(\alpha x+\beta y)^2$ schreiben, wobei $m,\alpha,\beta$ ganze Zahlen sind (Artikel 215). Setzen wir $\alpha x+\beta y=z$, so geht die vorgelegte Gleichung in diese über:
\[
mzz+2dx+2ey+f=0
\]
Daraus und aus $z=\alpha x+\beta y$ folgt
\[
x=\frac{\beta mzz+2ez+\beta f}{2\alpha e-2\beta d},\quad
y=\frac{\alpha mzz+2dz+\alpha f}{2\beta d-2\alpha e}
\]
Nun ist klar: Sofern nicht $\alpha e=\beta d$ gilt (diesen Fall betrachten wir gleich gesondert), genügen die Werte von $x$, $y$, die aus diesen Formeln durch Einsetzen eines beliebigen Werts für $z$ entstehen, der vorgelegten Gleichung. Es bleibt daher nur zu zeigen, wie sich die Werte von $z$ bestimmen lassen, die ganzzahlige Werte von $x$, $y$ liefern.

Wegen $\alpha x+\beta y=z$ kommen für $z$ notwendig nur \textit{ganze} Zahlen infrage. Macht außerdem ein Wert von $z$ sowohl $x$ als auch $y$ ganzzahlig, so liefern offenbar auch alle Werte von $z$, die zu ihm modulo $2\alpha e-2\beta d$ kongruent sind, ganzzahlige Werte. Setzt man also für $z$ alle ganzen Zahlen von $0$ bis einschließlich $2\alpha e-2\beta d-1$ (wenn $\alpha e-\beta d$ positiv ist) oder bis einschließlich $2\beta d-2\alpha e-1$ (wenn $\alpha e-\beta d$ negativ ist) ein, und werden für keinen dieser Werte sowohl $x$ als auch $y$ ganzzahlig, so liefert überhaupt kein Wert von $z$ ganzzahlige Werte von $x,y$, und die vorgelegte Gleichung ist nicht in ganzen Zahlen lösbar. Liefern dagegen einige jener Werte von $z$ ganzzahlige Werte von $x,y$, etwa $\zeta,\zeta',\zeta''$ usw. (die sich auch durch Lösen quadratischer Kongruenzen nach den Grundsätzen von Abschnitt IV finden lassen), so erhält man \textit{alle} Lösungen\SourceLeaf{225}, indem man $z=(2\alpha e-2\beta d)v+\zeta$, $z=(2\alpha e-2\beta d)v+\zeta'$ usw. setzt, wobei $v$ alle ganzen Zahlen durchläuft.

\Needspace{8\baselineskip}
\Article{220}
Für den ausgeschlossenen Fall, dass $\alpha e=\beta d$ gilt, müssen wir ein eigenes Verfahren finden. Nehmen wir an, dass $\alpha$, $\beta$ teilerfremd sind; nach Artikel 215, I ist das zulässig. Dann ist $\frac{d}{\alpha}=\frac{e}{\beta}$ eine ganze Zahl (Artikel 19), die wir $=h$ setzen. Die vorgelegte Gleichung nimmt dann diese Form an:
\[
(m\alpha x+m\beta y+h)^2-hh+mf=0
\]
Sie ist also offenbar nur dann rational lösbar, wenn $hh-mf$ eine Quadratzahl ist. Sei $hh-mf=kk$; dann sind die folgenden beiden Gleichungen offenbar der vorgelegten Gleichung äquivalent:
\[
m\alpha x+m\beta y+h+k=0,\quad\text{und}\quad
m\alpha x+m\beta y+h-k=0
\]
\textit{d.h.}, jede Lösung der vorgelegten Gleichung genügt auch einer dieser beiden Gleichungen und umgekehrt. Die erste Gleichung ist offenbar nur dann in ganzen Zahlen lösbar, wenn $h+k$ durch $m$ teilbar ist; ebenso lässt die zweite nur dann eine ganzzahlige Lösung zu, wenn $h-k$ durch $m$ teilbar ist. Diese Bedingungen reichen aber auch für die Lösbarkeit der jeweiligen Gleichung aus (da $\alpha$, $\beta$ als teilerfremd vorausgesetzt sind), und sämtliche Lösungen lassen sich nach den bekannten Regeln angeben.

\Needspace{8\baselineskip}
\Article{221}
Den in Artikel 217 betrachteten Fall (da er von allen der schwierigste ist) erläutern wir an einem Beispiel. Gegeben sei die Gleichung $xx+8xy+yy+2x-4y+1=0$. Daraus leiten wir zunächst durch Einführung der anderen Unbekannten $p=15x-9$. $q=15y+6$ die Gleichung $pp+8pq+qq=-540$ ab. Man findet, dass sämtliche ganzzahligen Lösungen dieser Gleichung in den folgenden vier Formeln enthalten sind:
\[
\begin{aligned}
p&=6t,&q&=-24t-90u\\
p&=6t,&q&=-24t+90u\\
p&=-6t,&q&=24t-90u\\
p&=-6t,&q&=24t+90u
\end{aligned}
\]
Dabei durchlaufen $t$, $u$ alle positiven ganzen Zahlen, die der Gleichung
\SourceLeaf{226}
$tt-15uu=1$ genügen. Sie sind in der Formel
\[
\begin{aligned}
t&=\tfrac{1}{2}\bigl((4+\sqrt{15})^n+(4-\sqrt{15})^n\bigr)\\
u&=\frac{1}{2\sqrt{15}}\bigl((4+\sqrt{15})^n-(4-\sqrt{15})^n\bigr)
\end{aligned}
\]
enthalten, wenn $n$ alle positiven ganzen Zahlen (einschließlich der Null) durchläuft. Daher sind sämtliche Werte von $x,y$ in den folgenden Formeln enthalten:
\[
\begin{aligned}
x&=\tfrac{1}{5}(2t+3),&y&=-\tfrac{1}{5}(8t+30u+2)\\
x&=\tfrac{1}{5}(2t+3),&y&=-\tfrac{1}{5}(8t-30u+2)\\
x&=\tfrac{1}{5}(-2t+3),&y&=\tfrac{1}{5}(8t-30u-2)\\
x&=\tfrac{1}{5}(-2t+3),&y&=\tfrac{1}{5}(8t+30u-2)
\end{aligned}
\]
Wendet man unsere Regeln richtig an, so findet man: Damit sich \textit{ganzzahlige} Werte ergeben, muss man in der ersten und zweiten Formel diejenigen Werte von $t,u$ nehmen, die aus einem \textit{geraden} Index $n$ entstehen; in der dritten und vierten dagegen diejenigen, die man aus einem \textit{ungeraden} $n$ erhält. --- Die einfachsten Lösungen sind $x=1,-1,-1$; $y=-2,0,12$, jeweils einander entsprechend.

Übrigens lässt sich die Lösung des in den vorhergehenden Artikeln erläuterten Problems meist durch verschiedene Kunstgriffe abkürzen, besonders beim Ausschließen unbrauchbarer Lösungen, \textit{d.h.} solcher, die Brüche enthalten. Doch um nicht zu ausführlich zu werden, haben wir dies hier übergehen müssen.
